\documentclass{UAmathtalk}

\author{Pallavi Dani}
\address{Louisiana State University}
\urladdr{https://www.math.lsu.edu/~pdani/}
\title{Large-Scale Geometry of Right-Angled~Coxeter Groups}
\date{Thursday, April 30, 2026}

\renewcommand*{\where}{Massry B012}

\begin{document}

\maketitle

\begin{abstract}
Right-angled Coxeter groups form an extremely accessible, yet remarkably rich class of objects in geometric group theory. They are defined by simple presentations: they are generated by involutions, with the only additional relations requiring certain pairs of generators to commute. Despite this elementary definition, they display an extraordinary range of geometric behaviors. Consequently, they have played a crucial role in the field, as a source of illuminating examples and counterexamples and as a testing ground for conjectures. In this talk, I will survey recent progress in understanding their large-scale geometry, focusing in particular on questions about quasi-isometry. Along the way, I will illustrate some of the main tools and techniques used for establishing such results, many of which are applicable in more general settings.
\end{abstract}

\end{document}
